% Suggested front matter matching the current manuscript wording.
% This file is not included by the build; papers/main.tex is authoritative.

% Suggested title
\title{Template Priors in Small CNNs: Activation-Patching Effects Depend on the Number of Patched Channels}

% Suggested abstract
\begin{abstract}
First-layer convolutional kernels can remain close to hand-designed templates
after training, but kernel-template similarity does not determine how internal
activation interventions affect the model output. We study this distinction in
two controlled rendering tasks and two small CNNs using template
initialization, continued template retention, and gradual release of the
retention penalty. Persistent retention strongly increases kernel-template
similarity. The release-versus-retention difference in selected-channel
fidelity, however, varies with the number of patched channels. In a prospective
confirmation on 20 previously unused \texttt{two\_concepts} TinyCNN blocks,
the predeclared contrast comparing larger ($k=4,8$) with smaller ($k=1,2$)
patch sizes is $B=0.33098$ (95\% CI $[0.27368,0.38827]$,
$p=2.28\times10^{-10}$). Release is strongly worse at $k=1,2$, slightly worse
at $k=4$, and slightly better at $k=8$. A separate 1,200-model study over 50
new blocks gives positive values of the same contrast for TinyCNN on both
tasks and across several prior schedules, whereas TwoLayerCNN gives different,
often negative estimates. An earlier positive four-channel selected-versus-
random fidelity difference is partly explained by a change in the random-
channel control, showing that the baseline can affect a relative patching
comparison. An independent 32-checkpoint reimplementation audit passes all
fixed tolerances. Within the tested models, kernel-template similarity and
activation-patching behavior therefore provide different information, and the
measured release-versus-retention effect depends on patch size and architecture.
The experiments do not establish a unique mechanism, human interpretability,
or natural-image generalization.
\end{abstract}

% Suggested contribution paragraph
The study combines a controlled comparison of initialization, continued
retention, and gradual release with a prospective test of patch-size dependence
and a larger secondary study across tasks, architectures, schedules, and all
channel counts. The independent checkpoint audit provides a separate
implementation check on the central stored measurements. The claims are
restricted to the tested synthetic renderer, first-layer interventions, and
small CNN family.
